\section{($\heartsuit$)複素積分の計算問題}

\subsection{Laurent展開}
留数解析を行うとき, 微分を使ってもいいのだが, なかなか計算が難しい場合はLaurent展開で考えることも有効である.

\begin{egprob} (RIMS, H24基礎)
$\alpha$を正の定数とするとき, 次の等式を示せ.
\[\int_{-\infty}^{\infty} \dfrac{e^{i\alpha x}e^x}{(1+e^x)^2}dx = \dfrac{2\alpha \pi}{e^{\alpha \pi} - e^{-\alpha\pi}}\]
\end{egprob}

\begin{screen}
\begin{point}
\tf{$e^x$を含む積分では, 高さ$2\pi$の横長の長方形経路が有効な場合があります.} これは$\exp$の周期性を利用する方法です. この経路では上辺と下辺の積分の関係が分かりやすく, 側辺の積分は十分遠方で0に収束することが期待できます.
\end{point}
\end{screen}

\begin{egans}
上のように経路を取る. つまり, $\pm R$, $\pm R + 2\pi i$を頂点とする長方形経路を反時計回りに積分する. 側面の積分は$R\to \infty$で0に収束する. 上辺の積分が下辺の積分の$-e^{-\alpha \pi}$倍になることも分かる. したがって, この経路に留数定理を適用すれば積分を計算できる. $e^{x} = -1 \iff z \in \pi i (2\mb{Z} + 1)$なので, 経路内の極は$i\pi$のみである. また, 2位の極であることも分かる. 微分による留数の計算は煩雑なので, $\dfrac{e^{i\alpha z}e^z}{(1+e^z)^2}$のLaurent展開を考える. $w = z-i\pi$とおく. $w\to 0$で,
\bealn{
\dfrac{e^{i\alpha z}e^z}{(1+e^z)^2} &= \dfrac{e^{i\alpha (w+i\pi)}}{e^z + e^{-z} + 2} \\
&= e^{-\alpha\pi}\dfrac{e^{iw}}{2 - e^w - e^{-w}} \\
&= e^{-\alpha \pi}\dfrac{e^w}{-w^2 + O(w^4)} \\
&= -e^{-\alpha\pi}\dfrac{e^w}{w^2}(1+O(w^2))\\
&= -e^{\alpha\pi}(\dfrac{1}{w^2} + \dfrac{1}{w} + O(1))(1+O(w^2)) \\
&= -e^{-\alpha\pi} (\dfrac{1}{w^2} + \dfrac{1}{w} + O(1))
}
よって留数は$-e^{-\alpha\pi}$であると分かる. これらから等式を示すことができる. \qed
\end{egans}
\begin{prac}
上の問の細部を埋めよ.
\end{prac}

\subsection{テクニック・公式まとめ}




\tbox{
\begin{theo}
「留数の公式」と同じ状況で, $\disp\int_{\del D}f(z)\, dz $は次で計算される。
\[2\pi i\disp\sum_{i:\, a_i\in A} \dfrac{1}{(m_i-1)!}\disp\lim_{z\to a_i}\dfrac{d^{m_i-1}}{dz^{m_i-1}}\left\{(z-a_i)^{m_i}f(z)\right\}\]
\end{theo}
}{title=留数定理}

\begin{rem}
{\Large \bf
$2\pi i$は決して忘れてはならない。実関数の積分なのに虚数が残ったり不自然に$\pi$が残っていたりしたら見返すこと。
}
\end{rem}

Laurent展開から留数を求めることも多い。$n$位の極を持っているときの留数の計算公式を以下に述べておく。
\begin{prop}{({\bf 留数の公式})}\\
$D$を領域とする。$f(z)$の(高々可算個の)異なる極の集合を$A=\{a_1,a_2, \cdots\}\neq \varnothing$とし, $D\setminus{A}$上で$f$は正則とする。\footnote{$A\neq\varnothing$としたので, 一致の定理より$A$は集積点を含まない。}$f(z)$が$z=a_i$において$m_i$の極を持つならば, 次が成立する。
\[\underset{z=z_0}{\mr{Res}} f(z) = \dfrac{1}{(m_i-1)!}\disp\lim_{z\to z_0}\dfrac{d^{m_i-1}}{dz^{m_i-1}}\left\{(z-z_0)^{m_i}f(z)\right\}\]
\end{prop}
以下の問題の状況はしばしば起こる。
\begin{egprob} (複素解析概論, 演習3.24, 87p)\\
$P,Q$が多項式で, $f(z) = e^{iz}P(z)/Q(z)$のとき, 次数について$\mr{deg} Q \geq \mr{deg} P + 1$ならば, $\disp\lim_{R\to \infty} \disp\int_{C_R}f(z) \, dz = 0$を示せ。ただし, $C_r: z=Re^{i\theta}$\, ($0\leq \theta\leq \pi$)である。これを使って, $\disp\int_{0}^{\infty} \dfrac{x\sin{x}}{x^2+1}\,dx$ を求めよ。
$$
\left(
\begin{tabular}{l}
{\bf Hint:} $0\leq \theta \leq \dfrac{\pi}{2}$において, 不等式$\sin{\theta}\geq \dfrac{2}{\pi}\theta$が成り立つ\\
ことに注意せよ。これは, {\bf Jordanの不等式}と呼ばれる。
\end{tabular}
\right)
$$
\end{egprob}





\subsection{留数解析}
留数解析系の複素積分は必ずできてほしい。
一度はやっておくとよい問題である。
\begin{egprob}
$\Delta$内にすべての零点をもつ$n$次多項式$(n\geq 2)$
\[P(z) = z^n + a_1z^{n-1} + \cdots + a_n \quad (a_i \in\mb{C})\]
を考える。
\ben
\item $\int_{\del D} \dfrac{dz}{P(z)}$を求めよ。
\item $\int_{\del D}\dfrac{z^ndz}{P(z)}$を求めよ。
\een
\end{egprob}
\begin{screen}
まさかすべての$\Delta$内の留数を足さなければならないのか？\tf{いや、そうではない！}$\infty$を含めたすべての留数和は$0$であるという事実を用いる！つまり, リーマン球面$\mb{CP}^1$の$w=x^{-1}$平面側の話にいれかえるのである。
\end{screen}


次のような応用もできる\footnote{Liが講究でやったやつをコピーした。}。
\begin{egprob} (\tf{余接関数の部分分数分解})
\begin{equation} \label{equation:1.1}
\dfrac{1}{\tau} + \sum_{d=1}^{\infty} \parena{\dfrac{1}{\tau - d} + \dfrac{1}{\tau + d}} = \pi \cot{\pi\tau} = \pi i - 2\pi i \sum_{m=0}^{\infty} q^m
\end{equation}
\end{egprob}
\prf
\step{1}{1番目の等号}
まず左辺が$\mb{C}\setminus{\mb{Z}}$で広義一様収束することを示す. $R>0$を固定し, 閉集合$|z| \leq R$で考える.$N\geq 2R$なる整数$N$をひとつ決める. 左辺を
\[\parena{\dfrac{1}{\tau} + \sum_{n=1}^{N-1} \dfrac{2z}{z^2-n^2} }+ \sum_{n=N}^{\infty} \dfrac{2z}{z^2-n^2}\]
と書くと, 第1項は有理関数であり, 収束性に関係しない.  第2項は, $n\geq N$ならば
\[\abs{\dfrac{2z}{z^2-n^2}} \leq \dfrac{2|z|}{n}\cdot \dfrac{1}{1-\dfrac{|z|^2}{n^2}} \leq \dfrac{2R}{n^2}\cdot \dfrac{1}{1- \dfrac{1}{4}} = \dfrac{8R}{3n^2}\]
だから, 第2項は$|z|\leq R$で一様収束する. よって左辺は$\{ |z| \leq R \}\setminus{\mb{Z}}$で正則である. $R$は任意だから$\mb{C}\setminus{\mb{Z}}$で広義一様収束する. \par
$N\in \mb{N}$(十分大), $R=N+\dfrac{1}{2}$とし, $\pm R \pm Ri$を頂点とする正方形の反時計回りの経路を$C$とする. $C$の内側に点$z\notin \mb{Z}$を取り, $|z| < R$であるようにする. 有理型関数$f$を
\[f(\zeta) = \dfrac{\pi \cot{\pi\zeta}}{\zeta-z}\]
を積分する. 留数は$\mb{Z}$の点$k$と$z$のみで
\[\res{\zeta = z}{f(\zeta)d\zeta} = \pi \cot{\pi z},\quad \res{\zeta=k}f(\zeta)d\zeta = \dfrac{1}{k-z} \]
である(Check!). よって, 留数定理により
\begin{eqnarray*}
\int_{C}f(\zeta) \dfrac{d\zeta}{2\pi i} &=& \pi\cot{\pi z} + \sum_{k=-N}^{N} \dfrac{1}{k-z}\\
&=& \pi\cot{\pi z} - \dfrac{1}{z} - \sum_{k=1}^{N} \parena{\dfrac{1}{z-k} + \dfrac{1}{z+k}}
\end{eqnarray*}
である.
\begin{claim}
$N\to \infty$のとき, $\disp\int_{C} f(\zeta) \dfrac{d\zeta}{2\pi i} \to 0$である.
\end{claim}
\prf  これは素朴にやると示せない. 補助的に$\dfrac{\pi\cot{\pi\zeta}}{\zeta}$を考える. これは偶関数であるため, 長方形の向かい合う辺の経路による積分は相殺されることに注意. よって
\begin{eqnarray*}
\int_{C}f(\zeta) \dfrac{d\zeta}{2\pi i} &=& \int_{C} \pi\cot{\pi \zeta}\parena{\dfrac{1}{\zeta - z} - \dfrac{1}{\zeta}}\, \dfrac{d\zeta}{2\pi i}\\
&=& z \int_{C} \dfrac{\pi\cot{\pi \zeta}}{\zeta(\zeta - z)} \, \dfrac{d\zeta}{2\pi i}
\end{eqnarray*}
$C$上$ |\cot{\pi z}| < 2$に注意せよ(Check!). よって, 上のノルムを評価すると
\[|z| \int_{C} \dfrac{\pi |\cot{\pi \zeta}| }{|\zeta|(|\zeta| - |z|)}\dfrac{|d\zeta|}{2\pi} \leq |z| \int_{C}\dfrac{|d\zeta|}{R(R-|z|)} = \dfrac{8R|z|}{R(R-|z|)}\]
だから, この積分は$N\to 0$で$0$に収束する. よって示された.\qed \par
複素積分が見えない問題にも, 複素積分は有効である. \tf{$[0,2\pi]$における積分は, $|z| = 1$での線積分と解釈できることがある. }
\begin{egprob}
\ben
\item $n$は自然数とするとき, $\disp\int_{0}^{2\pi} \dfrac{\sin{nx}}{\sin{x}}\, dx$を求めよ. (RIMS, H25基礎)
\een
\end{egprob}
\begin{egans}
(1): $t=e^{x}$のとき,
\[\dfrac{\sin{nx}}{\sin{x}} = \dfrac{t^n - t^{-n}}{t-t^{-1}} = \sum_{k=0}^{n-1} t^{k}t^{-(n-1-k)}\]
である. $\int_{0}^{2\pi} e^{imx}\, dx = 2\pi \delta_{0,m}$なのでここから計算できる.


\end{egans}



\begin{egprob}
$L_R (R>0)$は$-R+2i$を始点, $R+2i$を終点とする線分を表す. このとき次の値を求めよ.
\[\lim_{R\to \infty} \int_{L_R} \dfrac{\cos{z}}{z^2 + 1}\, dz\]
\end{egprob}



\begin{egprob} (2021基礎数学試験-6)
\[\int_{0}^{\infty} \dfrac{dx}{x^3-i}\]
\end{egprob}
\begin{egans}
$C_1 = \set{Ri\to 0}$とする.
\[\int_{C_1}\dfrac{dz}{z^3 - 1} = \int_{R}^{0} \dfrac{idt}{-it^3 - 1} = \int_{0}^{R} \dfrac{dt}{t^3 - i} \]
$C_2 = [0,R]$とすれば, 4分円の経路を考え, 円弧上の積分を無視し...


\end{egans}




\subsection{分枝を取る系}
\begin{egprob}
\[\int_{0}^{\infty} \dfrac{1}{}\]
\end{egprob}








\subsection{演習問題}

\subsubsection{Level A}
\begin{prob} (H22数学I)
次の積分を求めよ.
\[\int_{-\infty}^{\infty} \dfrac{e^{ix} - 1}{x(x^2+1)}dx\]
\end{prob}


\subsubsection{Level B}

\begin{egprob} ({\bf 複素積分集 その1})
　
\end{egprob}

\begin{egprob} ({\bf 複素積分集 その2})\footnote{2019年前期「複素函数論」(吉川氏)のレポ-ト問題から引用.}\\
$a$は正の実数とする。(5)では$m\in \mb{N}, -1 < \alpha < m-1$とする。以下の複素積分を計算せよ。
\beqn
\begin{array}{cccc}
(1)&　\disp\int_{0}^{\infty}\dfrac{dx}{1+x^6} & (2)& \disp\int_{-\infty}^{\infty} \dfrac{\cos{x}}{(x^2+1)^2} \,dx\\
(3)&　\disp\int_{0}^{\infty} \dfrac{\sin{x}}{x(x^2+a^2)}\,dx &(4)& \disp\int_{-\infty}^{\infty}\dfrac{\cos{(\pi x)}}{1+x^2+x^4} dx\\
(5)&　 \disp\int_{0}^{\infty} \dfrac{x^{\alpha}}{1+x^{m}} &(6)&　 \disp\lim_{r\to \infty} \disp\int_{-r}^{r}\dfrac{x\sin{(2\pi x)}}{1+x+x^2} \\
(7)& \disp\int_{-\infty}^{\infty} \dfrac{x\sin{x}}{(x^2+a^2)^2} \, dx&(8)&　 \disp\int_{0}^{\infty} \dfrac{(\sin{x})^3}{x^3(x^2+1)}\, dx \\
(9)& \disp\int_{-\infty}^{\infty} \dfrac{e^{(1+i)x}}{(e^{x} + 1)^2}\, dx&(10)&　 \disp\int_{0}^{\infty} \dfrac{1-\cos{x}}{x^2(x^2+a^2)}\, dx
\end{array}
\eeqn
\end{egprob}

\begin{prob} (東北選択H28) \label{prob:tohoku_H28_7}
\ben
\item $n\in \mb{N}$, $c\in \mb{C}$のとき$\int_{|z| = 1} \parena{z+\dfrac{c}{z}}^{n} \dfrac{1}{z}\, dz$を計算せよ。経路は反時計回り。
\item $m\in \mb{N}$, $a\in \mb{R}$のとき
{\scriptsize
\[\sum_{k=0}^{m}  \dfrac{(m!)^2(1+a)^{2k}(i-ai)^{2(m-k)}}{(2k)!(2m-2k)!}\int_{0}^{2\pi}\ (\cos{\theta})^{2k} (\sin{\theta})^{2m-2k}\, d\theta = 2\pi a^{m}\]
}
を示せ。
\een
\end{prob}


\subsubsection{Level C}


\subsubsection{Hints・Answer}
{\small
\begin{probans}
典型解法である. 上半円. $(1-\frac{1}{e})\pi i$.
\end{probans}

\begin{probans} %東北H28_7
(2): 一見すると扱いにくい問題です。偶数番目だけを取り出していることから, ある積分の実部または虚部を取ったものと考えられます。誘導に従って$I_n$を用います。ただし, 指数関数を三角関数に直すタイミングに注意してください。
\end{probans}
}

\newpage



\subsection{演習問題}
